\documentclass[11pt]{article}
\usepackage{amsmath,amssymb,amsthm,booktabs}
\usepackage[margin=1in]{geometry}
\newtheorem{theorem}{Theorem}
\newtheorem{lemma}[theorem]{Lemma}
\newtheorem{definition}[theorem]{Definition}
\title{Problem 879: Touch-screen Password}
\author{Project Euler Solutions}
\date{}
\begin{document}
\maketitle

\section{Problem Statement}

A touch-screen password is a sequence of $\ge 2$ distinct spots on an $n \times n$ grid. Rules:
\begin{enumerate}
\item Trace straight lines between consecutive target spots.
\item Unvisited intermediate spots on the line are automatically included.
\item Visited spots are removed from future consideration.
\end{enumerate}
Given $389{,}488$ passwords on $3 \times 3$, find the count for $4 \times 4$.

\section{Mathematical Foundation}

\begin{definition}[Grid and Collinearity]
Spots at $(i,j)$, $0 \le i,j \le n-1$. Two spots $A=(x_1,y_1)$, $B=(x_2,y_2)$ have intermediate lattice points iff $g = \gcd(|x_2-x_1|,|y_2-y_1|) > 1$.
\end{definition}

\begin{lemma}[Intermediate Spots]
The lattice points strictly between $A$ and $B$ are
\[
\left(x_1 + \frac{k(x_2-x_1)}{g},\; y_1 + \frac{k(y_2-y_1)}{g}\right), \quad k = 1, \ldots, g-1.
\]
\end{lemma}

\begin{proof}
Write $x_2-x_1 = gd_x$, $y_2-y_1 = gd_y$ with $\gcd(d_x,d_y)=1$. A point at parameter $t \in (0,1)$ is a lattice point iff $tgd_x, tgd_y \in \mathbb{Z}$. Since $\gcd(d_x,d_y)=1$, this requires $tg \in \mathbb{Z}$, giving $t = k/g$ for $k \in \{1,\ldots,g-1\}$.
\end{proof}

\begin{theorem}[Symmetry Group]
The $n \times n$ grid has dihedral symmetry $D_4$ of order~8. For $n=4$, the 16 spots form 3 orbits:
\begin{itemize}
\item 4 corners,
\item 8 edge-centers,
\item 4 inner spots.
\end{itemize}
\end{theorem}

\begin{proof}
Corners map to corners, edge non-corners to edge non-corners, and inner spots to inner spots under all 8 symmetries.
\end{proof}

\begin{lemma}[State Space]
DFS state $= (\text{current position}, \text{visited bitmask})$. Total states: $n^2 \cdot 2^{n^2}$.
\end{lemma}

\begin{proof}
$n^2$ positions, $2^{n^2}$ visited subsets. The pair determines all available moves.
\end{proof}

\section{Algorithm}

Enumerate via DFS with bitmask memoization:
\begin{enumerate}
\item Precompute intermediate spots for all pairs.
\item For each starting spot $s$, run DFS counting all paths of length $\ge 2$.
\item At each step, check intermediate spot requirements before moving.
\item Exploit $D_4$ symmetry to reduce starting positions to 3 orbit representatives.
\end{enumerate}

\section{Complexity Analysis}

\begin{itemize}
\item \textbf{Time:} $O(n^4 \cdot 2^{n^2})$. For $n=4$: $16^2 \cdot 2^{16} \approx 1.7 \times 10^7$ operations.
\item \textbf{Space:} $O(n^2 \cdot 2^{n^2})$ for memoization. For $n=4$: $\approx 10^6$ entries.
\end{itemize}

\section{Answer}

\[
\boxed{4350069824940}
\]

\end{document}
